% Part: computability
% Chapter: computability-theory
% Section: no-universal-function

\documentclass[../../../include/open-logic-section]{subfiles}

\begin{document}

\olfileid[fr]{cmp}{thy}{nou}
\olsection{Absence de fonction calculable totale universelle}

Il existe une fonction partielle calculable
universelle pour les fonctions partielles
calculables. En revanche, aucune fonction
\emph{totale} calculable n'est universelle
pour les fonctions calculables totales.

\begin{thm}
\ollabel{thm:no-univ}
Il n'existe aucune fonction calculable
\emph{totale} universelle pour les fonctions
calculables totales. Autrement dit, aucune
fonction totale et calculable
$\fn{Un}'(k,x)$ ne peut posséder la propriété
suivante : pour toute fonction calculable
totale~$f(x)$, il existe un entier naturel~$k$
tel que $f(x)=\fn{Un}'(k,x)$ pour tout~$x$.\footnote{Corrections éditoriales : l'introduction source écrit « total for the partial computable functions » là où le contexte exige « universal ». La reformulation du théorème omet ensuite l'hypothèse de totalité de $\fn{Un}'$ ; sans elle, la fonction universelle partielle précédente serait un contre-exemple.}
\end{thm}

\begin{proof}
La preuve est un argument diagonal simple.
Si $\fn{Un}'(k,x)$ était totale et calculable,
la fonction
\[
d(x) = \fn{Un}'(x, x) + 1
\]
serait elle aussi totale et calculable.
Or, par définition, $d(k)$ diffère de
$\fn{Un}'(k,k)$. Pour tout~$k$, les fonctions
$d(x)$ et $\fn{Un}'(k,x)$ diffèrent donc au
moins en un argument, à savoir $x=k$.
\end{proof}

\begin{explain}
Le
\olref[uni]{thm:univ-comp} montre comment
éviter cet argument diagonal : la fonction
universelle doit pouvoir être partielle.
Ainsi, $\fn{Un}$ est universelle même pour
les fonctions calculables totales, mais elle
n'est pas elle-même totale. L'argument
diagonal précédent ne s'applique donc pas
de la même manière.
\end{explain}

\begin{prob}
  Pour comprendre pourquoi l'argument diagonal
  de la preuve du
  \olref{thm:no-univ} ne s'applique pas dans
  le cas partiel, considérons la fonction
  $f(x) \simeq \fn{Un}(x,x)+1$.
  Est-elle partielle calculable ? Si oui,
  elle possède un indice~$e$, c'est-à-dire
  que $f(x) \simeq \fn{Un}(e,x)$.
  Que peut-on dire de~$f(e)$ ?
\end{prob}

\end{document}
